\documentclass[10pt,a4paper]{amsart}
\usepackage[margin=19mm]{geometry}
\usepackage{amsmath,amssymb,mathtools}
\usepackage[scr=rsfs,cal=euler]{mathalfa}
\usepackage{xfrac}
\usepackage[unicode,hidelinks,bookmarks=false]{hyperref}
\newcommand{\ie}{i.e.\ }
\newcommand{\cf}{cf.\ }
\newcommand{\resp}{respectively\ }
\newcommand{\Subsection}{\subsection}
\providecommand{\nobreakdash}{\nobreak\hbox{-}}
\setlength{\emergencystretch}{2em}
\multlinegap0pt
\usepackage{amssymb}
\usepackage{tikz}
\usepackage[shortlabels]{enumitem}
\usepackage{xcolor}
\newtheorem{theorem}{Theorem}
\newcommand{\Di}{\,\mathrm{d}}
\newcommand{\IP}{\mathbb{P}}
\newcommand{\NB}{\mathsf{NB}}
\newcommand{\dto}{\stackrel{d}{\longrightarrow}}
\newcommand{\floor}[1]{\lfloor #1 \rfloor}
\title{On Cycles in Multiset Permutations, Parking Functions, and Related Structures}
\author{Calum Buchanan, Fabian Burghart, Stephan Wagner,, Mei Yin}
\date{Source: arXiv:2606.19500v1 (CC BY 4.0)}
\begin{document}
\maketitle

\noindent\emph{Source and licence.} On Cycles in Multiset Permutations, Parking Functions, and Related Structures. Version arXiv:2606.19500v1 (CC BY 4.0), distributed by arXiv under the Creative Commons Attribution 4.0 International licence, http://creativecommons.org/licenses/by/4.0/, as recorded in the arXiv licence metadata for this exact version and shown on the abstract page https://arxiv.org/abs/2606.19500v1. Full licence notice: \texttt{licenses/arxiv-2606.19500-CC-BY-4.0.txt}.\par\smallskip

\noindent\textit{Editorial scope.} Counted unit: the article's theorem thm:regM (source0.tex lines 783--793). The article's complete original proof of it is delivered with it (source0.tex lines 795--825, one \texttt{proof} environment of 31 lines). No mathematics is written, repaired or re-derived: the statement and the proof are the article's own lines, in the article's own notation, and no retained line is substituted or re-flowed.  Retained uncounted as the source-local context the proof needs: lines 688--688.  Nothing was omitted from any retained span, so the package carries no omission record.  No retained line contains a citation, so the wrapper carries no bibliography and no bibliography entry.  No retained line contains a figure environment, an image-inclusion command or drawing code, so no image is delivered and the package declares no asset dependency.  Editorial formatting only, in addition: an AMS \texttt{amsart} preamble that restates the article's own declarations and macros used by the retained lines, namely the environments \texttt{theorem} on separate counters, together with the standard packages those lines need.  The retained environments print in this document as Theorem 1 in order of appearance, whereas the article prints the same items with the numbering of its own full document.  The complete native source of the article as distributed in the arXiv e-print for this exact version is bundled unchanged as \texttt{sources/203149/source0.tex}.
\subsection{Connected parking functions}\label{subsec:cpf}

\begin{theorem}\label{thm:regM}
 For $r\geq 2$ fixed and $M$ as above, we have
 \begin{equation}\label{eq:thmregM1}
  c_M(j)
  \sim j^{1/2} \frac{(jr-1)!}{r!^{j}} \sqrt{\frac{\pi r(r-1)}{2}}
 \end{equation}
 for $j\to\infty$. Moreover, let $L_j$ denote the number of cyclic points in a uniform random connected multiset permutation of $M$. Then
 \begin{equation*}
  j^{-1/2} L_j \dto \left|N\left(0,\frac{r}{r-1}\right)\right|.
 \end{equation*}
\end{theorem}

\begin{proof}
We argue as in Section \ref{subsec:cpf}, just with a negative binomial distribution instead of Poisson. We have
\begin{align*}
    c_M(j) &= \sum_{k=1}^j a(M,k)
    = \sum_{k=1}^j \frac{r-1}{r!^{j}} \binom{j}{k} k! r^k (jr-k-1)! \\
    &= \frac{r-1}{r!^{j}} j! \sum_{k=1}^j  \frac{r^k (jr-k-1)!}{(j-k)!}
    = \frac{r-1}{r!^{j}} j! \sum_{h=0}^{j-1}  \frac{r^{j-h} (jr-j+h-1)!}{h!} \\
    &= \frac{(r-1)r^j j!(jr-j-1)!}{r!^{j}}  \sum_{h=0}^{j-1} \binom{jr-j+h-1}{h} r^{-h} \\
    &= \frac{r^{jr} j!(jr-j-1)!}{(r-1)^{jr-j-1}r!^{j}}  \sum_{h=0}^{j-1} \binom{jr-j+h-1}{h} (1-r^{-1})^{jr-j} r^{-h} \\
    &= \frac{r^{jr} j!(jr-j-1)!}{(r-1)^{jr-j-1}r!^{j}} 
    \IP\big(\NB(jr-j,1-r^{-1})\leq j-1\big).
\end{align*}
The probability $\IP(\cdot)$ tends to $1/2$ as in the previous instances. Applying Stirling's formula, one gets
\begin{align*}
 \frac{j!(jr-j-1)!}{(jr-1)!} 
 &\sim \sqrt{\frac{2\pi j(jr-j-1)}{jr-1}} \cdot \frac{j^{j+jr-j-1}}{j^{jr-1}} \frac{\left(r-1-\frac{1}{j}\right)^{jr-j-1}}{\left(r-\frac{1}{j}\right)^{jr-1}}\\
 &\sim j^{1/2} \sqrt{\frac{2\pi (r-1)}{r}} \frac{(r-1)^{jr-j-1} e^{-1}}{r^{jr-1} e^{-1}}.
\end{align*}
Plugging this into the full expression for $c_M(j)$ leads to some more cancellations and \eqref{eq:thmregM1}.

For the second part, we again fix $x>0$, set $R=\floor{x\sqrt{j}}$ and compute
 \begin{align*}
  \IP( L_j \leq R)
  &= \frac{1}{c_M(j)} \sum_{k=1}^R a(M,k)
   \sim
    2\IP\big(j-R\leq \NB(jr-j,1-r^{-1})\leq j-1\big), \\
  &\sim 2\int_{-\frac{R}{\sqrt{j\frac{r}{r-1}}}}^{-\frac{1}{\sqrt{j\frac{r}{r-1}}}} \frac{1}{\sqrt{2\pi}} e^{-y^2/2} \Di y
   \sim \int_{-x}^x \frac{1}{\sqrt{2\pi \frac{r}{r-1}}} e^{-\frac{y^2}{2\frac{r}{r-1}}} \Di y,
 \end{align*}
 which is the cumulative distribution function of $\left|N\left(0,\frac{r}{r-1}\right)\right|$.
\end{proof}

\end{document}
